
Four experiments are designed to address the four research questions in sequence: from confirming existence (RQ1), to validating the measurement (RQ2), to revealing conditional structure (RQ3), to examining RLHF moderation (RQ4).

\subsubsection{4.1 Datasets}

\begin{table*}[htbp]
\centering
\resizebox{\dimexpr 2\columnwidth+\columnsep\relax}{!}{%
\begin{tabular}{llllll}
\toprule
Split & Task & Construction Method & n (adv) & n (clean) & Answer Options \\
\midrule
AdvGLUE MNLI & NLI (entailment) & Human-adversarial & 121 & 200 & 3 \\
ANLI R1 & NLI (entailment) & Model-in-loop (easy) & 200 & 200 & 3 \\
ANLI R2 & NLI (entailment) & Model-in-loop (medium) & 200 & 200 & 3 \\
ANLI R3 & NLI (entailment) & Model-in-loop (hard) & 200 & 200 & 3 \\
AdvGLUE QQP & Paraphrase (binary) & Human-adversarial & 78 & 200 & 2 \\
\bottomrule
\end{tabular}
}
\end{table*}

AdvGLUE \citep{Wang2021AdvGLUE} and ANLI \citep{Nie2020Adversarial} are selected because both preserve ground-truth labels by construction---a prerequisite for valid $\Delta ECE$ measurement. The three ANLI rounds provide a difficulty gradient within the model-in-loop construction method. AdvGLUE QQP provides a binary classification comparison with AdvGLUE MNLI to test task-type dependency within the same adversarial construction method.

\subsubsection{4.2 Baselines}

\textbf{RLHF comparison (RQ4):}
\begin{itemize}
\item \textit{Llama-2-7b-hf (base):} Pre-RLHF checkpoint; no instruction fine-tuning. Reference for adversarial calibration without alignment.
\item \textit{Llama-2-7b-chat:} RLHF-aligned checkpoint trained from the same 7b-parameter base. Comparison isolates RLHF fine-tuning at fixed model scale.
\end{itemize}

Temperature scaling was specified in the original experimental plan but was not applied in this pilot; it is identified as a planned extension.

\subsubsection{4.3 Evaluation Metrics}

\begin{itemize}
\item \textbf{$\Delta ECE$:} Primary outcome. Positive = degradation; negative = improvement.
\item \textbf{Label preservation rate:} For RQ2; must be $\geq 0.80$ for $\Delta ECE$ to be interpretable.
\item \textbf{$\Delta \Delta ECE$:} For RQ4; $\Delta ECE(base) - \Delta ECE(chat)$.
\item \textbf{RLHF moderation rate:} Fraction of adversarial cells where $\Delta \Delta ECE$ > 0.01 (base shows meaningfully more calibration degradation than chat). Gate criterion: $\geq 60$\% of cells.
\item \textbf{Accuracy ($\Delta Acc)$:} Reported alongside $\Delta ECE$ for context.
\end{itemize}

